
Repository maintenance prediction research has employed survival analysis, composite metrics, and supervised classification. He et al. (2024) predicted repository lifespan using gradient boosting with HITS centrality features on 103,354 GitHub repositories, achieving C-Index 0.810. Their approach required graph construction from user-repository star networks and distributed infrastructure. While effective, they did not establish whether simpler methods without graph features could achieve comparable accuracy.

Adejumo and Johnson (2025) proposed a Composite Stability Index (CSI) aggregating repository metrics with fixed weights (30\% activity, 25\% commits, 25\% issues, 20\% age), achieving F1 0.80 on 100 repositories. Their weighted-sum approach relied on manual weight tuning and did not compare against learned classifiers. Kuruppu and De Silva (2026) introduced an entropy-weighted index using six GitHub features (stars, forks, contributors, issues, commits, lines of code), demonstrating ''strong discriminative capability'' via tertile grouping and statistical tests, but did not report binary classification accuracy.

König et al. (2025) used gradient boosting for fault prediction on 2.4 million commits from 33 open-source projects, analyzing process metrics (churn, file age, revision frequency). Their work focused on code-level fault prediction rather than repository-level maintenance status. Li et al. (2026) demonstrated large-scale GitHub metadata extraction for 116,211 repositories using the REST API, establishing feasibility of data collection at scale but not addressing classification.

No prior work established a simple logistic regression baseline before testing complex methods. The gap addressed by the present study is the absence of controlled comparison between simple classification (logistic regression), manual aggregation (CSI-style metrics), and complex ensembles (gradient boosting) on the same dataset with the same features.

The present study is closest to He et al. (2024) in using supervised learning for maintenance prediction, but tests logistic regression first rather than deploying gradient boosting without baseline comparison. It is closest to Adejumo and Johnson (2025) in recognizing that repository metadata contains predictive signal, but uses learned classification rather than hand-tuned aggregation. The methodological contribution is testing the simplicity hypothesis explicitly: can logistic regression achieve $\geq 75$\% accuracy with six GitHub metadata features?
